\documentclass[journal]{IEEEtran}

\usepackage{amsmath,amssymb}
\usepackage{graphicx}
\usepackage{booktabs}
\usepackage{multirow}
\usepackage{todonotes}
\usepackage{cite}

\begin{document}

\title{FedQPNT: Federated, Quantum-Sensor-Augmented Positioning, Navigation and Timing with Continuous Attack-Adaptive Trust Fusion}

\author{\IEEEauthorblockN{Author names withheld for review}
\IEEEauthorblockA{Affiliation withheld for review}}

\maketitle

\begin{abstract}
Global Navigation Satellite System (GNSS) spoofing and jamming remain an open threat to autonomous
vehicles and unmanned aerial vehicles (UAVs), and existing defenses trade off three properties that
have not, to our knowledge, been combined in one closed loop: (i) federated learning (FL) across a
fleet so that one node's experience protects nodes that have not yet seen a given attack,
(ii) continuously varying, per-sensor trust rather than a binary detect-and-exclude decision, and
(iii) a physically independent quantum inertial sensor -- a cold-atom interferometer (CAI) -- whose
absolute, low-bias output is not subject to the same spoofing surface as GNSS. \todo{result: one-sentence
summary of the headline RMSE\_h(P\_att) comparison, once measured} We present FedQPNT, a system that
fuses a classical inertial measurement unit (IMU), GNSS, and a cold-atom accelerometer inside a
loosely-coupled error-state Kalman filter, in which a learnable, federated detector drives a continuous,
asymmetric trust law with formally bounded chattering, and in which detector parameters (never raw
sensor data) are aggregated fleet-wide by a trust-aware robust aggregator (TRIM-NB-R). The quantum
sensor's noise and fault statistics are calibrated against real cold-atom interferometer laboratory data
(d'Armagnac de Castanet et al., 2024) rather than assumed; the platform dynamics, GNSS observables, and
attacks are simulated at the raw-observable level. We report the system design, the pre-registered
hypotheses (H1--H4) that isolate the contribution of FL, continuous trust, and quantum aiding
respectively, and the experimental protocol. \todo{result: RMSE\_h(P\_att) FedQPNT vs B-cont}
\todo{result: RMSE\_h(P\_att) FedQPNT vs Baseline A (FL detect-and-exclude)}
\todo{result: latency\_on FedQPNT vs B-cont, cold-start scenario S8}
\todo{result: H3, latency\_eff with vs without quantum aiding, gradual spoofing}
We disclose two scoped limitations up front: the rigid-mode cold-atom sensor loses contrast above
$\Omega_c \approx 17\,\mathrm{mrad/s}$, making quantum aiding intermittent during ordinary vehicle
manoeuvres, and the loosely-coupled filter requires a GNSS covariance inflation factor $\kappa_R$ whose
root cause is still under investigation. No experimental outcome is claimed in this draft; all
quantitative statements are placeholders pending the completed evaluation matrix.
\end{abstract}

\begin{IEEEkeywords}
Federated learning, GNSS spoofing, cold-atom interferometry, quantum inertial sensing, sensor fusion, trust-aware fault tolerance, positioning navigation and timing (PNT).
\end{IEEEkeywords}

%=====================================================================
\section{Introduction}
%=====================================================================

\subsection{Motivation}
Positioning, navigation and timing (PNT) for autonomous vehicles and UAVs depends heavily on GNSS, a
signal that is inexpensive to spoof or jam and whose loss or manipulation directly threatens safety.
A wide literature addresses parts of this problem: federated learning has been used to detect GNSS
spoofing across a fleet without centralizing raw sensor data \cite{khan2025enhancing,chai2025navigation,liu2026g2fl,namagembe2026machine,zhou2023spatiotemporal};
classical and learned adaptive filters re-weight or exclude GNSS when it disagrees with inertial
propagation \cite{mehra1970identification,carlson1988federated,carlson1990federated,ren2020adaptive,pardhasaradhi2022gps,negru2024resilient};
and cold-atom interferometers (CAIs) have been proposed and increasingly demonstrated as inertial
sensors whose bias is orders of magnitude below that of classical accelerometers
\cite{wright2022cold,darmagnac2024atom}. Separately, quantum-enhanced and quantum federated learning
have been studied as a way to speed up or secure the \emph{learning computation} itself
\cite{hazarika2025quantum,chehimi2023foundations,chehimi2022quantum}. To date, however, these threads
have not been combined: no verified prior work couples (a) a physical quantum inertial sensor, (b)
fleet-wide federated learning of an attack detector, and (c) a continuous, attack-adaptive trust law
that keeps this detector's output as a smoothly varying confidence rather than a switch. The closest
robust-FL system for GNSS spoofing, G\textsuperscript{2}FL \cite{liu2026g2fl}, addresses robust
aggregation under Byzantine clients but does not couple its detector to a continuously varying,
per-sensor trust weight, nor does it use a physical quantum sensor.

\subsection{Gap statement}
\label{sec:gap}
No verified reference in our bibliography (Section~\ref{sec:related}, Table~\ref{tab:comparison})
implements federated learning together with a physical quantum inertial sensor, and among the FL-based
GNSS-defense papers we could verify, all use a binary detect-and-exclude action rather than a
continuously varying trust weight. Conversely, the continuous/adaptive-trust literature we verified
either has no federated learning component (it is single-node) or, when described as "federated", uses
the classical multi-sensor "federated filter" sense of the term (Carlson-style master/local filter
fusion within one platform) rather than federated learning across distributed clients
\cite{negru2024resilient} -- a distinction we make precise in Section~\ref{sec:related} because the
term is a recurring source of confusion in this literature. The specific, scoped gap this paper
addresses is therefore: \emph{a closed loop in which an FL-trained, fleet-shared spoofing/jamming
detector drives a continuous, formally bounded, per-sensor trust weight that gates a quantum-aided
navigation filter, with the quantum sensor's own noise and fault behaviour calibrated to real
laboratory data rather than assumed.}

\subsection{Contributions}
\begin{itemize}
\item A system architecture in which a single 15-state loosely-coupled error-state Kalman filter is
  shared, unmodified, across every baseline and ablation, so that comparisons isolate the contribution
  of each mechanism (FL, continuous trust, quantum aiding) rather than confounding it with
  implementation differences (Section~\ref{sec:method}).
\item A continuous, asymmetric trust law with a proven chattering bound
  ($T_{\mathrm{cyc}} \geq 26.1\,\mathrm{s}$ between trust cycles for any input sequence, adversarial
  included) that couples a learned, federated detector to the GNSS and quantum measurement
  covariances via a bounded inflation $R_{\mathrm{eff}} = R / \max(w, w_{\min})$
  (Section~\ref{sec:method}).
\item TRIM-NB-R, a trust-aware robust federated aggregator that combines norm clipping, coordinate-wise
  trimmed-mean/median aggregation, and a cosine-similarity reputation score with quarantine, evaluated
  against FedAvg and FedProx references under sign-flip, Gaussian, label-flip and "A Little Is Enough"
  poisoning (Section~\ref{sec:method}).
\item A quantum inertial sensor model whose noise density, per-interrogation-time noise scaling, and
  bursty fault process are calibrated from real cold-atom interferometer shot data
  \cite{darmagnac2024atom} rather than from published summary statistics alone, together with an
  explicit, honestly reported operating limit (rigid-mode contrast loss above
  $\Omega_c \approx 17\,\mathrm{mrad/s}$) that makes quantum aiding intermittent during ordinary vehicle
  manoeuvres (Section~\ref{sec:limitations}).
\item A pre-registered, mechanism-class baseline taxonomy (A: FL detect-and-exclude; B-bin: single-node
  detect-and-switch; B-cont: single-node continuous trust; B$'$: classical innovation-based adaptive
  filter) and four falsifiable hypotheses H1--H4, reported whatever their outcome
  (Section~\ref{sec:setup}).
\end{itemize}

\subsection{Paper organisation}
Section~\ref{sec:related} surveys related work in five groups and positions FedQPNT against it.
Section~\ref{sec:threat} specifies the node and threat model. Section~\ref{sec:method} gives the
fusion, trust, and federation equations. Section~\ref{sec:setup} describes the experimental protocol.
Section~\ref{sec:results} and Section~\ref{sec:discussion} report results and limitations (currently
placeholders). Section~\ref{sec:conclusion} concludes, followed by a reproducibility statement.

%=====================================================================
\section{Related Work}
\label{sec:related}
%=====================================================================

\subsection{Federated learning for GNSS security and spoofing detection}
Several recent papers apply federated learning to GNSS spoofing or jamming detection so that a fleet of
vehicles can share a detector without centralizing raw trajectory data. Khan et al.
\cite{khan2025enhancing} have each vehicle localize itself from dead-reckoning-style signals (yaw
rate, steering angle, wheel speed), compare this against GNSS, and use the resulting localization
mismatch to compute weights that are aggregated at a roadside unit, where a global support-vector
machine classifier is trained by FL; they report 99\% accuracy on CARLA-simulated trajectories. Chai et
al. \cite{chai2025navigation} similarly use federated learning to identify GNSS spoofing and jamming
signals for UAVs. Namagembe et al. \cite{namagembe2026machine} propose a machine-learning-based GPS
spoofing detection and mitigation scheme for UAVs. Zhou et al. \cite{zhou2023spatiotemporal} propose a
spatial-temporal federated transfer learning architecture with multi-sensor data fusion for cooperative
positioning across vehicles; this targets positioning accuracy under heterogeneous data distributions
rather than adversarial spoofing specifically, but it establishes that federated transfer learning is a
workable mechanism for multi-vehicle positioning problems. Closest to our own robust-aggregation
component is G\textsuperscript{2}FL \cite{liu2026g2fl}, which proposes a robust federated learning
scheme explicitly for GNSS spoofing detection under Byzantine or poisoned clients. In every one of these
papers, the FL-trained detector's output drives a \emph{binary} decision at the node
(classify-and-exclude, or classify-and-flag); none couples the detector to a continuously varying
trust weight on the fusion filter's measurement covariance, and none uses a physical quantum sensor.
TRIM-NB-R differs from G\textsuperscript{2}FL's robustness mechanism in aggregation detail (per-round
cosine-similarity reputation with quarantine, rather than G\textsuperscript{2}FL's own scheme), but the
central point of departure from this entire group is architectural: in FedQPNT the FL-trained detector
is one input to a continuous trust law, not the exclude decision itself.

\subsection{Adaptive and fault-tolerant sensor fusion and trust}
The idea of re-weighting or excluding a sensor based on its own innovation statistics predates federated
learning by decades. Mehra \cite{mehra1970identification} introduced innovation-based adaptive
estimation of measurement and process noise for Kalman filters, the root of the "adaptive Kalman filter"
family and the basis of our B$'$ baseline. Carlson's federated filter
\cite{carlson1988federated,carlson1990federated} decomposes a navigation filter into local filters per
sensor feeding a master filter, giving fault tolerance and decorrelation properties; this is the
classical, single-platform sense of "federated," and it is worth stating precisely because the term
recurs, confusingly, in more recent papers. Ren et al. \cite{ren2020adaptive} demonstrate adaptive
sensor fusion of camera, GNSS and IMU for autonomous driving, adjusting fusion weights as sensor quality
changes; we use this as our representative of the single-node, continuous-weighting mechanism class
(B-cont in our taxonomy) because, unlike \cite{mehra1970identification}, it targets exactly the
camera/GNSS/IMU setting our own filter generalizes. Pardhasaradhi and Cenkeramaddi
\cite{pardhasaradhi2022gps} detect GPS spoofing on a drone via distributed radar tracking and, once
spoofing is flagged, switch control to the independent radar track; this is a detect-then-switch
mechanism against an independent second sensor, which we use as the representative of our B-bin
(detect-and-switch) baseline class, while noting that it switches to a different physical sensor
(radar) rather than continuously re-weighting the same GNSS/INS filter. Negru et al.
\cite{negru2024resilient} integrate GNSS, IMU, a monocular camera and a barometer using local
extended Kalman filters feeding a gated-recurrent-unit master filter; despite being described in its own
title as a "hybrid federated fusion architecture," this is the classical Carlson-style federated
\emph{filter} architecture with one learned component inside a single UAV, not federated learning across
distributed clients, and we cite it here, in the fusion-architecture group, rather than alongside the
FL-for-spoofing-detection papers of Section~2.1 -- a reframing that sharpens the gap statement of
Section~\ref{sec:gap}. A companion paper from the same group \cite{negru2024uwb} studies UWB-aided
hybrid navigation in degraded GNSS environments. None of this group's continuous or adaptive fusion
mechanisms is trained or shared across a fleet by federated learning, and none uses a physical quantum
sensor as an independent aiding source.

\subsection{Quantum inertial sensing for navigation}
Cold-atom interferometers measure acceleration or rotation via atom-interferometric phase, giving an
absolute output whose bias is set by fundamental atomic constants rather than by the drift of a
mechanical or MEMS proof mass. Wright et al. \cite{wright2022cold} review cold-atom inertial sensors for
navigation applications, covering the physical principle, achievable sensitivities, and the practical
challenges (dead time, dynamic range, size/weight/power) of fielding such a sensor on a moving platform.
d'Armagnac de Castanet et al. \cite{darmagnac2024atom} report a three-axis hybrid cold-\textsuperscript{87}Rb
atom interferometer operated at arbitrary orientations and rotation rates, providing, to our knowledge,
the most directly usable real dataset for calibrating a moving-platform CAI noise and fault model: per-shot
Raman population-ratio residuals at several interrogation times, and a measured contrast-versus-rotation-rate
law in both a rigid-mirror and an actively tip-tilt-compensated ("inertial-pointing") mode. We use this
dataset directly (Section~\ref{sec:threat}, and D-014/D-015/D-019/D-020/D-021 in the project decision
log) to calibrate our simulated CAI's noise density, its per-interrogation-time scaling, and a
Gilbert--Elliott bursty-outlier channel, rather than assuming these from the summary numbers in the
paper's text alone. Neither \cite{wright2022cold} nor \cite{darmagnac2024atom} addresses fleet-wide
learning or GNSS attack robustness; they are the physical-sensing foundation our system builds on, not a
navigation-security contribution in themselves.

\subsection{Quantum federated learning (distinguished from quantum sensing)}
A separate line of work applies quantum computation to the federated learning process itself, and it is
important to distinguish it from quantum \emph{sensing}, since the two are sometimes conflated. Hazarika
et al. \cite{hazarika2025quantum} propose quantum-enhanced federated learning for metaverse-empowered
vehicular networks, using quantum-enhanced computation to accelerate or improve FL convergence, not a
physical quantum sensor providing navigation-grade measurements. Chehimi et al.
\cite{chehimi2023foundations} lay foundational architectures for quantum federated learning over
classical and quantum communication networks, and a related paper by the same first two authors
\cite{chehimi2022quantum} studies quantum federated learning specifically with quantum \emph{data}
(i.e., data that is itself a quantum state, transmitted over a quantum network) -- again, a statement
about the learning substrate and communication channel, not about inertial or positioning sensors. None
of this group's methods involves a physical quantum sensor of the kind used in FedQPNT (a cold-atom
accelerometer providing a classical-unit acceleration measurement), and none targets GNSS spoofing or
jamming. We cite this group specifically to pre-empt a natural but mistaken reading of "quantum
federated learning for navigation" as prior work equivalent to ours: it is not, since it addresses the
learning computation rather than the sensing modality.

\subsection{GNSS spoofing and jamming signatures}
Detecting GNSS attacks depends on the observable signatures those attacks leave, at the level of raw
receiver measurements rather than final position. Psiaki and Humphreys \cite{psiaki2016gnss} give a
foundational survey of GNSS spoofing mechanisms and detection principles, including the
capture-and-carry-off dynamics that a power-advantage spoofer must use to smoothly pull a receiver's
tracking loops away from the true signal, and the vulnerability of a receiver at the moment of
re-acquisition after an outage. An earlier study by Humphreys et al.
\cite{humphreys2008assessing} assesses the practical spoofing threat and its feasibility with
civilian-grade equipment. Radoš et al. \cite{rados2024recent} survey recent advances in jamming and
spoofing detection in GNSS and report field-measured signatures we rely on directly, including the
collapse of cross-satellite carrier-to-noise-density-ratio (C/N0) correlation under single-antenna
spoofing (from $-0.76$ clean to $0.99$ spoofed) -- a signature our attack and detector-feature models
reproduce explicitly, since a single spoofing antenna necessarily transmits all counterfeit signals with
a common-mode power fluctuation. Borhani-Darian et al. \cite{borhanidarian2024detecting} detect GNSS
spoofing with a deep-learning classifier operating on the cross-ambiguity function in the correlation
domain, a complementary, lower (pre-tracking) level of the receiver signal chain than the
raw-observable (pseudorange/Doppler/C-N0/AGC) level at which our own attacks and detector features
operate. Tariq and Ahanger \cite{tariq2026multimodel} propose a multi-model fusion scheme -- Kalman
filtering combined with ensemble machine learning and a transformer-based detector -- for detecting and
mitigating BeiDou spoofing across a decentralized UAV swarm, reporting approximately 99\% detection
accuracy and a greater-than-97\% mission-success rate under active spoofing; we note, following our own
verification, that this fusion is per-model ensembling rather than an explicit decentralized-consensus
mechanism, and that the frequently quoted "97\%" figure in this line of work is mission-success rate,
not detection accuracy. Together, this group establishes the observable-level attack signatures
(C/N0 correlation collapse, AGC and clock-jump behaviour, RAIM blindness to slow drift-off spoofing) that
motivate our own attack simulator and detector feature set (Section~\ref{sec:threat}).

\begin{table*}[t]
\centering
\caption{Comparison of related systems against the four properties FedQPNT combines. ``FL?'' = detector or model trained/aggregated across distributed clients (not the classical single-platform ``federated filter'' sense). ``Quantum?'' = uses a physical quantum sensor (not quantum-enhanced computation). ``Adaptive trust?'' = continuously varying, not binary, sensor trust during an active attack. ``Attack-tested?'' = evaluated under an explicit spoofing/jamming/poisoning scenario in the cited paper.}
\label{tab:comparison}
\begin{tabular}{lcccc}
\toprule
System / paper & FL? & Physical quantum sensor? & Adaptive (continuous) trust during attack? & Attack-tested? \\
\midrule
Khan et al. 2025 \cite{khan2025enhancing} & Yes & No & Not stated (detection only) & Yes (GPS spoofing) \\
Chai et al. 2025 \cite{chai2025navigation} & Yes & No & No (classification only) & Yes (spoofing/jamming) \\
G\textsuperscript{2}FL, Liu et al. 2026 \cite{liu2026g2fl} & Yes (robust) & No & No & Yes (spoofing, Byzantine clients) \\
Negru et al. 2024 \cite{negru2024resilient} & No (classical federated \emph{filter}) & No & Continuous (GRU master filter) & No (degraded GNSS, not adversarial) \\
Pardhasaradhi \& Cenkeramaddi 2022 \cite{pardhasaradhi2022gps} & No & No & No (detect-and-switch to radar) & Yes (GPS spoofing) \\
Hazarika et al. 2025 \cite{hazarika2025quantum} & Yes & No (quantum-enhanced compute) & No & No \\
Chehimi et al. 2023 \cite{chehimi2023foundations} & Yes (quantum FL) & No & No & No \\
Tariq \& Ahanger 2026 \cite{tariq2026multimodel} & No (multi-model fusion) & No & Partial (model ensembling) & Yes (BeiDou spoofing) \\
\textbf{FedQPNT (this paper)} & \textbf{Yes} & \textbf{Yes (cold-atom, real-data-calibrated)} & \textbf{Yes (continuous, bounded-chatter)} & \textbf{Yes (15-scenario matrix)} \\
\bottomrule
\end{tabular}
\end{table*}

%=====================================================================
\section{System and Threat Model}
\label{sec:threat}
%=====================================================================

\subsection{Node model}
The system is a fleet of $N$ independent nodes (ground vehicles or UAVs), each running the same code
path. Every node is split into two halves connected only by defined sensor-output types: an
\emph{Environment} half (ground truth trajectory, sensor models, GNSS signal generation, and any
attacks) and an \emph{Agent} half (GNSS receiver, fusion filter, trust engine, learnable detector, and
FL client). Only sensor outputs cross from Environment to Agent; ground truth (\texttt{TruthState}), the
per-tick attack label (\texttt{AttackLabel}), and the environment's own private metadata on
\texttt{GnssEpoch} never do, and this boundary is enforced both by construction (the environment strips
private fields before handing an epoch to the receiver) and by an automated static-analysis guard that
scans the Agent-side packages for forbidden imports of truth or label types. A node runs a 100~Hz global
tick; at each tick the IMU and, when scheduled, the quantum sensor and GNSS receiver are stepped, giving
a fixed per-tick sequence (truth $\rightarrow$ sensors $\rightarrow$ optional attack $\rightarrow$
receiver $\rightarrow$ filter propagation/innovations $\rightarrow$ feature extraction $\rightarrow$
detector $\rightarrow$ trust update $\rightarrow$ filter correction $\rightarrow$ pseudo-labelling and FL
buffering). Nodes exchange only federated-learning model updates with a server, on a fixed round
schedule (Section~\ref{sec:method}); navigation never blocks on the availability of a fresh global
model, so a node that has lost communications keeps navigating and detecting with its last-known
detector.

\subsection{Sensors}
Each node carries three sensor types feeding one shared, loosely-coupled 15-state error-state Kalman
filter (position, velocity, attitude error, accelerometer bias, gyroscope bias): a classical strapdown
IMU (accelerometer plus gyroscope, one of two grades -- MEMS or tactical -- reported separately because
their free-inertial error growth differs by roughly an order of magnitude); a simulated GNSS receiver
that computes a position/velocity/clock fix by weighted least squares with receiver autonomous integrity
monitoring (RAIM) from per-satellite pseudorange, pseudorange-rate, carrier-to-noise-density ratio and
automatic-gain-control (AGC) observables; and a cold-atom accelerometer (CAI) used in a
hybrid-with-classical configuration, where the CAI's absolute, low-bias output is fused as a direct,
dynamics-independent observation of the classical accelerometer's bias (both instruments sense the same
specific force over the same time window, so gravity, attitude and trajectory terms cancel in the
hybridisation residual). The CAI's noise density, per-interrogation-time scaling, and a bursty
outlier/fringe-jump fault process are calibrated from real per-shot atom-interferometer data
\cite{darmagnac2024atom} rather than assumed from published summary statistics: field-grade operation
uses a $2T = 20\,\mathrm{ms}$ interrogation time, a robust per-shot noise of $\sigma_{\mathrm{shot}}
\approx 5.6\,\mu g$ (within 11\% of the value implied by the paper's own stated signal-to-noise ratio),
and a two-state Gilbert--Elliott channel fitted to a measured $9.0\%$ outlier rate with $66\%$ burst
persistence. The CAI keeps reporting \texttt{valid=True} during a burst -- it is silently wrong, exactly
as a real sensor would be -- so the trust engine, not a validity flag, is what must catch a bad quantum
sample (Section~\ref{sec:method}).

\subsection{Attacker capabilities and limits}
The attacker acts only on the GNSS signal, at the raw-observable level (per-satellite pseudorange,
pseudorange rate, C/N0, lock state, AGC/noise floor), never on the IMU, the quantum sensor, or the
inter-node communication channel in the scenarios reported here. Five attack families are modelled: (i)
gradual (``drift-in'') spoofing, in which a single-antenna spoofer captures the receiver's tracking loops
via a brief power-advantage bump and then drags the code and Doppler measurements smoothly toward a
fake trajectory, remaining statistically close to indistinguishable from clean GNSS to a snapshot RAIM
test; (ii) meaconing (record-and-rebroadcast), which appears as a common-mode clock/delay shift across
all satellites with only a small C/N0 rise; (iii) abrupt, non-aligned spoofing, which cannot capture
smoothly and instead forces a momentary loss of lock, a reacquisition delay, and a discontinuous
position jump on re-lock; (iv) continuous-wave and wideband jamming, which raises the receiver noise
floor and reduces the number of locked satellites; and (v) combined jam-then-spoof, in which a denial
phase is followed by a capture attempt at the moment of reacquisition -- the classic window in which a
receiver is most vulnerable to being handed to a spoofer \cite{psiaki2016gnss}. A single-antenna
spoofer's signature (common-mode C/N0 fluctuation across all spoofed satellites, and the loss of the
normal elevation dependence of C/N0) is modelled explicitly, following field-measured correlation
statistics reported by Radoš et al. \cite{rados2024recent}, and is available to any detector operating
at the raw-observable level, though it is not itself exposed as a hand-crafted detector feature by the
attack layer. In the federated-learning threat surface, a fraction $f$ of nodes may be Byzantine,
submitting arbitrary model updates (sign-flipped, Gaussian-noised, label-flipped, or the "A Little Is
Enough" family of small, aggregate-shifting perturbations) rather than honest gradients; the
self-reported sample count on an update is untrusted, since an attacker could claim an arbitrarily large
$n$ to bias a naive weighted average.

\subsection{What never leaves a node}
No raw sensor stream, trajectory, or attack label ever leaves a node. The only object that crosses a node
boundary to the server is a federated \texttt{ModelUpdate}: a detector weight delta (882 floating-point
parameters, about 3.5~kB), a coordinate-wise running mean/standard-deviation pair for feature
normalisation, and a small set of scalar training metrics (sample counts, loss, pseudo-label rate). The
server never receives ground truth, position, velocity, or raw features. Trust-law hyperparameters and
the fusion filter itself are not learned or shared in this version of the system; they are fixed,
identical across every method under comparison, and tuned once, before any test-seed result is examined,
on a disjoint range of tuning seeds (Section~\ref{sec:setup}).

%=====================================================================
\section{Method}
\label{sec:method}
%=====================================================================
\todo[inline]{This section is an equation-level outline; the surrounding exposition, worked examples and
figures are to be expanded in a later drafting pass.}

\subsection{Fusion filter and trust-to-covariance coupling}
The shared 15-state error-state Kalman filter processes GNSS position/velocity and the CAI-derived
hybridisation residual as independent innovations. For a sensor $s$ with continuous trust weight
$w_s \in [w_{\min}, 1]$, trust couples into the filter only through the effective measurement covariance
and an exclusion rule, so that every method (including a fixed-trust baseline) shares one fusion core:
\begin{align}
R_{\mathrm{eff},s} &= \frac{R_s}{\max(w_s,\, w_{\min})}, \label{eq:reff}\\
&\text{skip the update entirely if } w_s < w_{\mathrm{excl}}. \label{eq:excl}
\end{align}
A $\chi^2$ normalised-innovation-squared (NIS) gate, evaluated with $R_{\mathrm{eff}}$, is applied to
every sensor and every method alike, so a low-trust sensor can still contribute a small, gated update
rather than being invisibly discarded.

\subsection{Continuous, asymmetric trust law}
Trust is updated at each GNSS evidence event $j$ from a smoothed detector output $\bar p_j$ (an
exponentially smoothed version of the learned detector's spoof/jam probability) via a target trust
$\tau^*_j$ and an asymmetric, hysteresis-gated recursion toward it:
\begin{align}
\bar p_j &= \bar p_{j-1} + \beta_j\,(p_j - \bar p_{j-1}), && \beta_j = 1 - e^{-\Delta t_j/\tau_p}, \label{eq:smooth}\\
\tau^*_j &= w_{\min} + (1-w_{\min})\,\mathrm{clip}\!\left(\frac{\theta_{\mathrm{hi}} - \bar p_j}{\theta_{\mathrm{hi}} - \theta_{\mathrm{lo}}},\,0,\,1\right), \label{eq:target}\\
w_j &=
\begin{cases}
w_{j-1} + (1-e^{-\Delta t_j/\tau_d})(\tau^*_j - w_{j-1}) & \text{if } \tau^*_j < w_{j-1} \\
w_{j-1} + (1-e^{-\Delta t_j/\tau_r})(\tau^*_j - w_{j-1}) & \text{if } \tau^*_j > w_{j-1} \wedge G_j \\
w_{j-1} & \text{otherwise,}
\end{cases} \label{eq:trustlaw}
\end{align}
where $G_j$ is a recovery gate requiring both a sustained clean detector output and a run of
nominal-range NIS values before trust is allowed to recover, and the fast distrust time constant
$\tau_d$ is much smaller than the slow recovery time constant $\tau_r$. Equation~\eqref{eq:trustlaw} is a
convex combination of $w_{j-1}$ and $\tau^*_j \in [w_{\min}, 1]$ at every step, so $w$ is bounded in
$[w_{\min}, 1]$ by construction, and for constant $\tau^*$ the law is a contraction with exponential,
overshoot-free convergence. Counting a ``trust cycle'' as a downward crossing of $w=0.5$ followed by an
upward crossing of $w=0.9$, the recovery gate's minimum dwell time on the smoothed statistic implies a
provable lower bound on the cycle period, $T_{\mathrm{cyc}} \geq 26.1\,\mathrm{s}$, for any input
sequence, including an adversarial one -- this is what bounds trust chattering (evaluated in scenario
S7, Section~\ref{sec:setup}) independent of the detector's own quality.

\subsection{TRIM-NB-R: trust-aware robust federated aggregation}
Each round, live nodes submit a detector weight delta $\Delta_i$. The server aggregates as follows:
\begin{align}
\Delta_i &\leftarrow \Delta_i \cdot \min\!\left(1,\; \frac{c \cdot \mathrm{median}_j \lVert \Delta_j\rVert}{\lVert \Delta_i \rVert}\right), && c = 2, \label{eq:clip}\\
\Delta &= \mathrm{coordinate\text{-}wise\ trimmed\ mean}_{\beta=0.2}(\{\Delta_i\}_{i\ \mathrm{live}}), \label{eq:trim}\\
r_i &\leftarrow \rho\, r_i + (1-\rho)\max(0, \cos(\Delta_i, \Delta)), && \rho = 0.8, \label{eq:reputation}\\
\theta_g &\leftarrow \theta_g + \eta_s \Delta, && \eta_s = 1, \label{eq:serverstep}
\end{align}
with nodes whose reputation $r_i$ falls below a threshold for several consecutive rounds quarantined for
a fixed number of subsequent rounds. Norm clipping bounds any single node's influence before aggregation;
the trimmed mean (falling back to a coordinate-wise median when too few nodes are live to trim safely)
tolerates up to a designed fraction $\beta = 20\%$ of Byzantine nodes; the reputation and quarantine step
adds a slower, cross-round defense against nodes that are individually within norm but persistently
inconsistent with the consensus update direction. FedAvg and FedProx are retained as reference
aggregators (the former weights by each node's self-reported sample count, which is exactly the
untrusted quantity a poisoning node can misreport, and is expected to be the more vulnerable reference
for that reason).

\subsection{Local training and pseudo-labels}
No node ever has access to a ground-truth attack label. Each node trains its detector locally against
\emph{hindsight pseudo-labels}: an epoch is labelled positive if a sustained innovation-based CUSUM
statistic, computed against the CAI-aided filter as a slowly drifting but trustworthy hindsight
reference, crosses a threshold within a fixed look-ahead window, or if AGC, RAIM or clock-jump features
cross conservative, independently justified thresholds; it is labelled negative only if every feature
stays within its nominal range for the entire window; otherwise it is not used for training. This is the
specific mechanism by which the quantum sensor contributes to the learning problem, not only to
navigation: a CAI-aided inertial reference that drifts far more slowly than GNSS position error under
denial is a more trustworthy hindsight label source than GNSS alone, and federated learning is what lets
one node's hindsight-labelled experience of an attack family improve the shared detector's real-time,
causal decision at a node that has not yet experienced that family.

%=====================================================================
\section{Experimental Setup}
\label{sec:setup}
%=====================================================================
\todo[inline]{Expand with concrete parameter tables, hardware/runtime budget, and exact seed ranges once
the M1/M2 integration described in the project log completes.}

Every method under comparison -- FedQPNT and its baselines and ablations (A, A0, B-bin, B-cont, B$'$,
and the $-$quantum, $-$FL, fixed-trust, binary, $-$recovery-gate, logreg and FedAvg ablations) --
shares the same sensor noise realisations (paired, seed-indexed
random number generator streams), the same fusion filter and NIS gate, the same detector architecture and
initial weights, and the same tuning budget, so that a difference in outcome can be attributed to the
mechanism under test rather than to incidental implementation differences (a design constraint drawn
directly from the project's scientific-integrity rule that results must be reported exactly as measured,
with no scenario dropped after the fact). \todo[inline]{Table: baseline/ablation matrix (FedQPNT, A, A0,
B-bin, B-cont, B$'$, $-$quantum, $-$FL, fixed-trust, binary, $-$recovery-gate, logreg, FedAvg) with
detector-training / trust-law / CAI / aggregator columns.}

\subsection{Hypotheses}
Four pre-registered, falsifiable hypotheses are reported whatever their outcome:
H1 (FedQPNT improves on Baseline A, FL detect-and-exclude, on $\mathrm{RMSE}_h(P_{\mathrm{att}})$ and
$\mathrm{MAX}_h(P_{\mathrm{att}})$); H2 (FedQPNT improves on B-cont, the strongest single-node
continuous-trust competitor, on the same metrics and on detection latency for attack families a given
node has not yet experienced); H3 (FedQPNT with quantum aiding improves on the $-$quantum ablation on
detection latency for gradual spoofing); H4 (FedQPNT improves on B-cont on a cold-start node's detection
probability and latency). The pre-registered primary metrics are $\mathrm{RMSE}_h(P_{\mathrm{att}})$
(horizontal position RMSE during the active-attack phase) and $\mathrm{latency}_{\mathrm{on}}$
(time from attack onset to sustained detection); all other metrics are secondary.

\subsection{Scenario matrix}
Fifteen scenarios (nominal operation; gradual spoofing at three severities; sudden jamming; combined
jam-then-spoof at reacquisition; partial node failure/delayed FL updates; long-duration CAI bias drift;
trust chattering under toggled spoofing; cold-start node join; communications dropouts; sample-rate
mismatch; extreme sensor noise; trust-score/model poisoning at two Byzantine fractions; post-attack
recovery; multi-hour stability; and simultaneous attacks on a fleet subset) exercise the system along
axes independent of the scientific hypotheses, each with its own pre-registered, pass/fail engineering
acceptance criterion evaluated over 30 paired test seeds with a Wilcoxon signed-rank test as the primary
statistical test and Holm--Bonferroni correction across the confirmatory family
$\{H1,\ldots,H4\}\times\{\text{primary metrics}\}$.

%=====================================================================
\section{Results}
\label{sec:results}
%=====================================================================
\todo[inline]{No results exist yet. This section will report, per scenario and per method: RMSE\_h/MAX\_h/P95\_h and ANEES for S1--S15; detection probability and latency\_on/latency\_eff for the spoofing/jamming scenarios; trust chattering counts for S7; FL global-model AUC, rounds skipped and quarantined-node counts for S5/S8/S9/S12/S15; and the H1--H4 hypothesis tests with Hodges--Lehmann effect sizes and Holm-corrected p-values. All numbers below are placeholders.}

\todo{result: Table -- RMSE\_h(P\_att), MAX\_h(P\_att), P95\_h(P\_att) for FedQPNT, A, A0, B-bin, B-cont, B$'$, per attack scenario (S2--S4, S15), 30 seeds, median [IQR]}
\todo{result: Table -- detection probability P\_D and latency\_on/latency\_eff, same methods and scenarios}
\todo{result: Figure -- trust trajectory $w_{\mathrm{gnss}}(t)$, $w_{\mathrm{quantum}}(t)$ and detector output $\bar p(t)$ across a representative gradual-spoof run}
\todo{result: Table -- S12 poisoning: global-model AUC drop at $f=20\%$ and $f=40\%$, TRIM-NB-R vs FedAvg vs FedProx}
\todo{result: S1 nominal-operation cost check (RMSE\_h ratio to fixed-trust, FAR/h, ANEES, mean $w_{\mathrm{gnss}}$) for every method}
\todo{result: H1--H4 hypothesis test table with Hodges--Lehmann median differences, 95\% BCa CIs, Holm-adjusted p-values}
\todo{result: MEMS vs tactical IMU grade comparison of the CAI's contribution (H3), reported for both grades per D-002/D-009 risk R-6}

%=====================================================================
\section{Discussion and Limitations}
\label{sec:discussion}
\label{sec:limitations}
%=====================================================================

The following limitations are already established by the project's own validation work and are stated
here rather than deferred, in keeping with the scientific-integrity requirement that known weaknesses be
disclosed regardless of how they affect the headline comparison.

\begin{itemize}
\item \textbf{The quantum sensor's calibration is real, the platform dynamics are not.} The CAI noise
  density, per-interrogation-time scaling and bursty-outlier fault process are fitted to real cold-atom
  interferometer laboratory shot data \cite{darmagnac2024atom}, but the vehicle/UAV trajectories, the
  GNSS constellation and signal propagation, and the resulting sensor time series along those
  trajectories are all simulated; no public dataset places a moving quantum inertial sensor under GNSS
  attack, so a fully real-data closed-loop evaluation is not currently possible.
\item \textbf{Rigid-mode quantum aiding is intermittent during ordinary manoeuvres.} The measured
  rigid-mode contrast roll-off rate implies a rotation-rate tolerance of only
  $\Omega_c \approx 17\,\mathrm{mrad/s}$ (field grade), well below the rotation rates of an ordinary
  vehicle turn ($\sim 0.3\,\mathrm{rad/s}$). The CAI therefore loses contrast, and its aiding becomes
  intermittent, during manoeuvring rather than only during straight-and-level segments. This is reported
  as a real effect, not corrected for, and it is expected to weaken the measured quantum benefit relative
  to an idealized always-on quantum sensor.
\item \textbf{The GNSS is simulated at the raw-observable level, not the signal (IF-sample) level.}
  Pseudorange, Doppler, carrier-to-noise-density ratio, lock state and AGC are modelled directly, with
  attacks constructed as literature-grounded deltas on the clean observables; this keeps published
  attack signatures faithful and keeps multi-hour, multi-node simulation computationally tractable, but
  it cannot capture receiver-internal (tracking-loop, correlator) effects that an IF-sample-level or
  hardware-in-the-loop evaluation would.
\item \textbf{The loosely-coupled filter requires a GNSS covariance inflation factor $\kappa_R$.} Under
  fixed trust, the nominal-operation filter is over-confident unless the GNSS fix covariance is inflated
  by a currently non-unit factor; this is being investigated as an effect of the loosely-coupled fix's
  time-correlated (ionospheric/tropospheric/multipath) error being processed as if it were white, and,
  separately, of the receiver's reported velocity-fix covariance not yet being independently checked for
  honesty. $\kappa_R$ is tuned once, identically for every method, on tuning seeds disjoint from the test
  seeds, so it cannot bias a between-method comparison; its root cause is documented as an open,
  in-progress investigation rather than resolved.
\item \textbf{The novelty claimed here is scoped, not absolute.} G\textsuperscript{2}FL
  \cite{liu2026g2fl} already provides robust federated learning for GNSS spoofing detection under
  Byzantine clients. FedQPNT's contribution is the closed loop that couples an FL-trained detector to a
  continuous, formally bounded trust law and to physical quantum-sensor aiding, not the robust-aggregation
  idea in isolation.
\item \textbf{Hypotheses are pre-registered; no result yet exists.} H1--H4 above are stated as hypotheses,
  to be reported whatever their outcome, including a null or reversed result, per the project's own
  scientific-integrity rule.
\end{itemize}

%=====================================================================
\section{Conclusion}
\label{sec:conclusion}
%=====================================================================
\todo[inline]{To be finalized once Section~\ref{sec:results} is populated. Draft closing paragraph
below.}

We have presented FedQPNT, a system architecture that closes the loop between federated, fleet-shared
attack detection, a continuous and formally bounded per-sensor trust law, and physical quantum inertial
aiding whose noise and fault behaviour is calibrated against real laboratory data rather than assumed.
We have specified the node and threat model, the fusion and trust equations, the federated aggregation
mechanism, and a pre-registered, mechanism-class baseline and hypothesis structure designed so that the
eventual results can be attributed to a specific mechanism rather than to incidental implementation
choices. \todo{result: one-sentence outcome summary once the evaluation matrix completes} We have also
disclosed, in advance of any result, the two most consequential limitations of the current design: the
intermittency of rigid-mode quantum aiding above ordinary vehicle turn rates, and the as-yet-unresolved
need for a GNSS covariance inflation factor in the loosely-coupled filter.

\section*{Reproducibility Statement}
Every numeric default in the system specification is tagged as literature-derived, design-derived, or an
explicit assumption requiring a sensitivity sweep before publication (project decision log, D-001 and
architecture specification conventions). All Monte Carlo experiments use a hierarchical, deterministic
random-number-generator seeding scheme (one independent stream per node, component and purpose) so that
adding or removing a component does not change the noise realization seen by any other component, and so
that paired statistical tests across methods are exact rather than approximate. Pre-training, tuning and
test seed ranges are disjoint and fixed before any test-seed result is examined; the trust-law and filter
hyperparameters are frozen, with their configuration hash logged, before the test-seed sweep begins. The
quantum-sensor calibration pipeline, from the raw Zenodo-hosted dataset \cite{darmagnac2024dataset}
through per-interrogation-time robust noise statistics to the bursty-outlier channel, is scripted
end-to-end and is re-run by an automated test suite. No experimental result, ablation, or scenario is
withheld or dropped after being observed; any exclusion, if one becomes necessary, is logged with its
reason in the project decision log. Code, configuration hashes, and the full 15-scenario result set will
be released alongside the camera-ready version.

\bibliographystyle{IEEEtran}
\bibliography{refs}

\end{document}
